\chapter{Implementation}
\label{ch:impl}

\section{Introduction}
This chapter describes how the design of Chapter~\ref{ch:design} was built. It
begins with the technology stack and the reasoning behind each choice, then
follows the data as it moves through the system: ingestion, index computation,
feature engineering, the supervised builder, the three pipelines, the model
layer, the promotion gate and the registry, inference, explanation, simulation
and monitoring. It then covers the serving layer, the continuous-integration
workflows and the production deployment, and closes with the failures encountered
during development and what each of them required.

Code excerpts are presented as figures, captioned below and centred, in keeping
with the convention that the document numbers only figures and tables. Excerpts
are lightly abridged for length; the behaviour described is that of the running
system.

\section{Technology Stack}
Table~\ref{tab:stack} records the technology choices and their justification.

\begin{xltabular}{\textwidth}{P{2.6cm} P{3.1cm} Y}
\caption{Technology stack and the reason for each choice.}\label{tab:stack}\\
\toprule
\bfseries Layer & \bfseries Choice & \bfseries Justification \\
\midrule
\endfirsthead
\toprule
\bfseries Layer & \bfseries Choice & \bfseries Justification \\
\midrule
\endhead
\bottomrule
\endlastfoot
Language & Python 3.11 & The machine-learning ecosystem; pinned below 3.12 for
compatibility across the dependency set at the time of development. \\
Data source & Open-Meteo~\cite{ref:openmeteo} & Free, \emph{keyless}, hourly,
global, and serving both multi-year history and a multi-day forecast from the
same request shape. The absence of a credential removes an entire class of
operational failure from the scheduled jobs. \\
Feature store and model registry & Hopsworks Serverless~\cite{ref:hopsworks} & A
managed free tier providing both halves of the state plane, with primary-key
upserts, event time and metric-tagged model versions. \\
Orchestration & GitHub Actions~\cite{ref:ghactions} & Free scheduled compute,
co-located with the source, with secret injection and concurrency groups. \\
Tabular processing & pandas~\cite{ref:pandas} & The standard library for the
grouped, shifted, windowed operations that the feature design requires. \\
Classical models & scikit-learn~\cite{ref:sklearn} & Ridge, random forest,
metrics, standardisation and the time-series splitter used for the backtest. \\
Champion model & XGBoost~\cite{ref:xgboost} & The strongest candidate measured;
trains a 200,000-row, 27-feature problem in about a minute on CPU, which is what
makes a nightly retrain feasible. \\
Sequence-model study & TensorFlow / Keras~\cite{ref:tensorflow} & Used to build
and evaluate the LSTM alternative; deliberately excluded from the deployment
image. \\
Explainability & SHAP (TreeExplainer)~\cite{ref:treeshap} & Exact, polynomial-time
attributions for tree ensembles, with the additive property that makes
plain-language rendering possible. \\
Backend & FastAPI and Uvicorn~\cite{ref:fastapi} & Typed, asynchronous, with
request-model validation and automatic interface documentation. \\
Frontend & React with Vite~\cite{ref:react,ref:vite} & Component model suited to
a multi-tab dashboard; Vite produces a small static bundle deployable to an edge
host. \\
Charting and maps & Recharts and react-leaflet~\cite{ref:leaflet} & Responsive
charts with band series for the prediction interval, and an interactive map of
Pakistan. \\
Advisor & Gemini or Groq (free tiers), or Claude, over the Model Context
Protocol~\cite{ref:mcp} & Structured tool calling, so answers are grounded in the
system's own forecasts. The provider is chosen from whichever API key is present,
so the advisor needs no paid account; MCP publishes the same tools to external
clients. \\
Deployment & Docker on Railway; Vercel~\cite{ref:docker} & A container gives full
control of the backend image; the frontend needs no server at all. \\
Testing & pytest~\cite{ref:pytest} & Unit tests over the index computation, the
feature builder and the supervised assembly. \\
\end{xltabular}

\section{Configuration as a Single Source of Truth}
One module defines everything the rest of the system needs to agree on: the
project root and the derived data and model directories, the credentials read
from the environment, the feature-store object names and versions, the forecast
horizon, and the list of supported cities. Each city is an immutable record of
name, coordinates, province and timezone.

The consequence is that adding a city is a single line, and that no module
anywhere else in the codebase contains a hard-coded path, city name or object
version. The path derivation deserves note because it later mattered in
deployment: the project root is computed relative to the configuration module's
own location, which is why the deployment image installs the package in editable
mode rather than copying it into the interpreter's library directory - an
ordinary install would have resolved the root to a location containing none of
the committed artefacts.

\section{Data Ingestion}
Ingestion wraps three provider endpoints behind one helper. The air-quality
endpoint serves both history and forecast; weather is split between a forecast
endpoint and an archive endpoint derived from ERA5~\cite{ref:era5}, and the
helper selects between them according to whether an explicit date range was
requested. Seven pollutant variables and ten weather variables are fetched and
merged on the observation hour, and the merged frame is tagged with the city name
and its coordinates, giving 21 raw columns.

The session is configured to retry, because these jobs run unattended. A momentary
throttle or gateway error from the provider would otherwise fail an entire hourly
run; instead it is retried with exponential backoff on precisely the status codes
that indicate a transient condition, as Figure~\ref{fig:code-retry} shows.

\begin{figure}[htbp]\centering
\begin{lstlisting}[language=Python]
def _session(retries: int = 4, backoff: float = 1.5) -> requests.Session:
    """Retry transient failures: this runs unattended every hour."""
    retry = Retry(
        total=retries,
        backoff_factor=backoff,          # waits 1.5s, 3s, 6s, ...
        status_forcelist=(429, 500, 502, 503, 504),
        allowed_methods=frozenset(["GET"]),
    )
    session = requests.Session()
    session.mount("https://", HTTPAdapter(max_retries=retry))
    return session
\end{lstlisting}
\codegap
\caption{The retrying HTTP session used for every provider request.}
\label{fig:code-retry}
\end{figure}

\section{The Air Quality Index}
\label{sec:impl-aqi}
The index module implements the EPA definition~\cite{ref:epa} directly: the six
index bands, the concentration breakpoints per pollutant, the piecewise-linear
sub-index formula, the maximum-over-sub-indices rule, and the six health
categories with their advisory text. The categories are used throughout the
system - for the colour scale, the alert wording and the advisor's answers - so
defining them once here keeps the interface, the API and the language model in
agreement.

Two implementation decisions in this module matter.

The first is the treatment of gases. As noted in Chapter~\ref{ch:literature}, the
EPA specifies gaseous breakpoints in volumetric mixing ratios while the provider
reports mass concentrations. Rather than apply an approximate conversion, the
computed index is restricted to the two particulate species, whose breakpoints
are already expressed in the units the provider uses. This is a defensible
restriction on this data - particulate matter dominates the index in every city
studied - and it avoids a class of error that would be invisible in aggregate
metrics, since the index is a maximum and one wrong sub-index would silently
dominate every affected row. The restriction is asserted by a unit test that
feeds an enormous carbon monoxide value alongside clean particulate readings and
requires the index to be unaffected.

The second is precedence. Where the provider supplies its own index value it is
preferred, because the provider applies the correct EPA averaging windows; the
locally computed value is used only to fill gaps. The module thus serves both as
the fallback computation and as the documentation of what the predicted number
actually means.

\section{Feature Engineering}
The feature builder takes a raw hourly frame and produces the 74-column
engineered table specified in Chapter~\ref{ch:design}. It runs in a fixed order,
because the change, lag and rolling families all require the target column to
exist first.

The two properties that make the feature table trustworthy are visible in
Figure~\ref{fig:code-roll}. Every rolling window is shifted by one observation
before aggregation, so that it terminates at the previous hour, and every family
is computed inside a single city, because the multi-city entry point groups by
city and engineers each group independently before concatenating.

\begin{figure}[htbp]\centering
\begin{lstlisting}[language=Python]
def _add_rolling_features(df: pd.DataFrame) -> pd.DataFrame:
    # The .shift(1) is the anti-leakage guard: the window ends at the
    # PREVIOUS hour and never includes the current observation.
    for col in _SERIES_COLS:
        for win in _ROLL_WINDOWS:
            roll = df[col].shift(1).rolling(
                win, min_periods=max(2, win // 2))
            df[f"{col}_roll_mean_{win}h"] = roll.mean()
            df[f"{col}_roll_std_{win}h"]  = roll.std()
            df[f"{col}_roll_max_{win}h"]  = roll.max()
    return df


def build_features_all_cities(raw_all, *, dropna_target=True):
    # Lags must never reach across a city boundary: Lahore's value from
    # 24 hours ago is meaningless for Karachi.
    parts = [build_features(g, dropna_target=dropna_target)
             for _, g in raw_all.groupby("city", sort=False)]
    return pd.concat(parts, ignore_index=True)
\end{lstlisting}
\codegap
\caption{The two mechanisms that prevent leakage: a shifted rolling window and
per-city grouping.}
\label{fig:code-roll}
\end{figure}

The builder also computes the feature group's primary key, converting the
observation datetime to epoch seconds by subtracting the Unix epoch and dividing
by a one-second interval. This slightly indirect formulation was adopted because
the direct conversion helper's behaviour differs across major pandas versions,
and the primary key must be stable for the upsert semantics to hold.

\section{The Multi-Horizon Supervised Builder}
The supervised builder implements Section~\ref{sec:design-mh}. For each city it
sorts by time and then, for each of the ten horizons, constructs a block in which
the target-time features are taken from the current row, the anchor-state
features are taken from $h$ rows earlier by shifting, the horizon column is set to
$h$, and the target is the index at the current row. Blocks are concatenated,
rows with an incomplete anchor are dropped, and the result is subsampled with a
fixed seed. Figure~\ref{fig:code-sup} shows the core.

\begin{figure}[htbp]\centering
\begin{lstlisting}[language=Python]
for city, g in features.groupby("city", sort=False):
    g = g.sort_values("datetime").reset_index(drop=True)
    for h in horizons:
        block = pd.DataFrame()
        for col in TARGET_TIME_FEATURES:     # known for the future hour tau
            block[col] = g[col]
        for col in ANCHOR_STATE_FEATURES:    # observed h hours before tau
            block[f"{col}_anchor"] = g[col].shift(h)
        block["horizon"] = h                 # the horizon IS a feature
        block[TARGET_COLUMN] = g["aqi"]      # AQI at tau
        block["city"], block["datetime"] = city, g["datetime"]
        blocks.append(block)
\end{lstlisting}
\codegap
\caption{Assembly of the multi-horizon supervised set. One observation
contributes up to ten rows, differing only in the horizon and the target.}
\label{fig:code-sup}
\end{figure}

The chronological split is implemented as a quantile of the datetime column
rather than a positional cut, so that it remains correct even when the rows are
unevenly distributed in time after subsampling. The datetime column is coerced to
a datetime type inside the split function rather than being assumed, for a reason
discovered in production and discussed in Section~\ref{sec:impl-challenges}.

\section{The Feature Pipeline}
\label{sec:impl-feature-pipeline}
The hourly pipeline is deliberately short: it fetches the last seven days for all
22 cities, engineers features per city, and upserts the result. Its two
non-obvious properties are both correctness properties.

It fetches seven days rather than one hour, because the lag and rolling features
for the newest observations need history behind them; requesting only the newest
hour would produce a row whose 24-hour features were all missing. The resulting
overlap with data already stored is not a problem, because the write is an upsert
on the composite key: re-running the pipeline, or running it twice in the same
hour, cannot create a duplicate row. Idempotence of this kind is what allows the
recovery strategy for every failure in this system to be ``run it again''.

\section{The Backfill Pipeline}
The backfill runs the same feature logic over the full history. It is triggered
manually and it is the component that required the most hardening, because it is
the only job that runs for hours and touches every city.

It fetches each city's history in three-month chunks and concatenates them before
engineering features, so that lags and rolling windows are computed over a
contiguous series rather than being reset at every chunk boundary. It writes each
city as that city completes, rather than accumulating everything and writing once,
so that an interruption costs at most one city. It catches per-city exceptions and
continues, so that one failing city does not abort the other twenty-one.

Its resume logic is history-aware, and the distinction matters. An earlier version
skipped any city already present in the store; but the hourly pipeline seeds every
city with recent rows within an hour of being enabled, so on the next run every
city looked present and the backfill skipped all of them, silently leaving the
store with a week of history instead of three years. The corrected logic asks a
different question - does this city's earliest stored observation reach back to
the requested start date? - and only then treats the city as complete.
Figure~\ref{fig:code-resume} shows it.

\begin{figure}[htbp]\centering
\begin{lstlisting}[language=Python]
# History-aware resume: a city counts as done only if its data
# already reaches back to `start`. A city with merely recent rows
# (seeded by the hourly pipeline) is NOT done and gets backfilled.
coverage = city_coverage()          # {city: earliest date present}
done = {c for c, earliest in coverage.items()
        if earliest is not None and earliest <= start_date}
partial = {c for c in coverage if c not in done}
\end{lstlisting}
\codegap
\caption{History-aware resume in the backfill. Presence in the store is not
evidence of completeness.}
\label{fig:code-resume}
\end{figure}

\section{The Storage Layer}
The storage facade exposes four operations and selects a backend by whether a
credential is configured, degrading to the local columnar backend when the
managed client is not installed. That second condition is not hypothetical: on
the development platform the managed client cannot be installed at all, because
one of its transitive cryptographic dependencies fails to build there. The facade
turns a hard platform incompatibility into a configuration detail.

The managed backend concentrates every vendor-specific call in one module. Three
of its choices were forced by the free tier and are recorded here as findings.
The feature group is created with a time-travel table format that materialises
server-side, because the platform's newer default format requires a client
library that the base installation does not ship, and group creation fails
outright without it. Statistics computation is disabled, because the profiler over
74 columns exhausted the free executor's memory and failed the materialisation
job; the statistics are used only by the platform's own interface and by no
pipeline. And inserts do not block on materialisation, for the reason given in
Figure~\ref{fig:code-insert}.

\begin{figure}[htbp]\centering
\begin{lstlisting}[language=Python]
def insert_features(project, df, *, wait: bool = False) -> None:
    """Upsert features (non-blocking by default).

    On the free tier a Spark executor occasionally hangs at
    INITIALIZING; a blocking insert then freezes the whole backfill
    until the scheduler kills it at its time cap. Non-blocking is
    safe here because inserts are idempotent primary-key upserts
    and the backfill resumes by coverage.
    """
    fg = get_feature_group(project)
    fg.insert(_sanitize(df), write_options={"wait_for_job": wait})
\end{lstlisting}
\codegap
\caption{Non-blocking, idempotent feature insertion.}
\label{fig:code-insert}
\end{figure}

\section{The Training Pipeline}
The daily training pipeline implements use case UC-5. It reads the whole feature
table, builds the supervised set, splits it chronologically, scores the
persistence baseline by simply reading the anchor column, fits the three learned
candidates, and selects the challenger by validation RMSE.

The candidate configurations are fixed rather than searched, because a nightly
unattended job that performs a hyper-parameter search is both slower and less
reproducible than one that does not, and because the measured gap between the
families is far larger than the gap between plausible configurations within a
family. The gradient-boosted candidate uses 600 trees at a learning rate of 0.05,
a maximum depth of 8, and row and column subsampling at 0.8 with the histogram
tree method; the forest uses 120 trees at a maximum depth of 16 with a minimum
leaf size of 10; the linear candidate is a standardised Ridge regression. Every
candidate has a fixed random seed, satisfying NFR-7.

Prediction intervals are fitted immediately after selection, on the same
validation window, by grouping the residuals by horizon and taking the tenth and
ninetieth percentiles within each group, as Figure~\ref{fig:code-interval} shows.

\begin{figure}[htbp]\centering
\begin{lstlisting}[language=Python]
valid["pred"]  = best.predict(X_va)
valid["resid"] = valid[TARGET_COLUMN] - valid["pred"]
interval_by_horizon = {
    int(h): [float(g["resid"].quantile(0.1)),
             float(g["resid"].quantile(0.9))]
    for h, g in valid.groupby("horizon")
}
\end{lstlisting}
\codegap
\caption{Per-horizon empirical residual quantiles, giving an 80\% prediction
interval that widens with lead time.}
\label{fig:code-interval}
\end{figure}

Everything inference needs is then packaged into a single bundle: the fitted
estimator, the exact feature column order, the modelled horizons, the interval
table, the metrics of every candidate, the winning family's name and the training
timestamp. Shipping the column order with the estimator rather than relying on
both sides importing the same constant is a small but deliberate guard against
training-serving skew.

\section{The Evaluation Suite}
\label{sec:impl-evaluate}
Evaluation is implemented as a separate module rather than inside the training
pipeline, because it answers a different question. Training asks which candidate
is best; evaluation asks how good the promoted model actually is, and at what
lead time.

Per-horizon metrics are computed by predicting once over the whole validation set
and then masking by horizon, so that RMSE, MAE and $R^2$ are reported separately
for each of the ten modelled lead times. The persistence baseline is scored under
exactly the same mask, which is what makes the comparison meaningful: a claim to
beat a baseline at 72 hours is only informative if the baseline was also
evaluated at 72 hours.

The walk-forward backtest uses an expanding-window time-series split with five
folds. Each fold trains a fresh gradient-boosted model on the folds that precede
it and tests on the fold that follows, so that every test window lies strictly in
the future of its own training data. This is the construction that the
literature~\cite{ref:tashman,ref:bergmeir} identifies as the honest estimate of
forward performance, and it is deliberately more pessimistic than the single
chronological split used for candidate selection, because each early fold is
trained on far less data.

The module writes four artefacts: two figures, a markdown report, and a JSON
document that the API serves and the dashboard renders. Publishing the same
numbers into a document that reaches the public interface, rather than into a log
that only a developer reads, is what satisfies principle P6. The results
themselves are presented in Chapter~\ref{ch:testing}.

\section{The Champion-Challenger Gate}
\label{sec:impl-gate}
The gate is the component that turns a retraining script into a governed system.
Figure~\ref{fig:gate} illustrates it, and Figure~\ref{fig:code-gate} shows the
rule.

\fig{champion-challenger}{0.92}{The promotion gate. Both outcomes are recorded;
the leaderboard rows are the system's actual history.}{fig:gate}

\begin{figure}[htbp]\centering
\begin{lstlisting}[language=Python]
PROMOTION_MARGIN = 0.0   # 0 = any improvement promotes; a tie does NOT

champion   = current_champion(entries)
challenger = candidates[best_name]
version    = max((e.get("version", 0) for e in entries),
                 default=0) + 1

promoted = champion is None or \
           challenger["rmse"] < champion["rmse"] - PROMOTION_MARGIN
if promoted:
    for e in entries:            # demote the previous champion
        e["is_champion"] = False
\end{lstlisting}
\codegap
\caption{The promotion rule. Strict inequality means an exact tie keeps the
incumbent.}
\label{fig:code-gate}
\end{figure}

Three details are worth drawing out. The comparison is a strict inequality, so a
challenger that merely matches the champion is refused; a model does not ship
because it is newer. The first ever run is promoted unconditionally, because
there is no incumbent to beat. And every run appends an entry recording not only
the winner but the metrics of every candidate considered, the previous champion's
score and the decision - so the leaderboard is an audit trail rather than a
scoreboard.

The gate's behaviour in production, shown in the lower half of
Figure~\ref{fig:gate}, is the clearest demonstration that it works. Version~1 was
promoted at an RMSE of 20.597 as the first champion. Version~2, trained after the
full historical backfill had completed, reached 19.687 and displaced it.
Version~3, trained an hour later on effectively identical data, reached the same
19.687 - and was \emph{refused}, because 19.687 is not less than 19.687. That
refusal is the system declining to churn its own production model for no gain.

Only if the gate promotes is the bundle written, and the write goes to two places:
a local file for use within the same job, and the durable model registry. The
registry is what satisfies principle P1. A model saved only to the runner's disk
ceases to exist minutes later; a model registered with its metrics can be
retrieved by any later job, including a deployed backend that was never present
when training happened. Retrieval asks the registry for the version with the
lowest RMSE, which makes ``the champion'' a property of the stored metrics rather
than of a mutable pointer that could drift out of step.

\section{The Model Layer and the Sequence-Model Study}
The registry facade presents one pair of operations - save and load - and hides
which backend answers them, preferring the durable registry and falling back to
the local bundle. This is what allows the identical inference code to run on a
continuous-integration runner that has a registry credential and inside a
deployment container that deliberately has neither the credential nor the client.

Alongside the production model, a recurrent network was implemented to test the
assumption that a sequence model would do better. It consumes 48 hours of
pollution and weather history for a city as a sequence and predicts the index 24
hours ahead, through two stacked LSTM layers followed by a dense head with
dropout, trained with early stopping on a chronological validation split.
Sequences are built per city with a stride, windows containing missing values are
discarded, and scaling is fitted on the training portion only so that no
validation statistic leaks backwards.

The comparison was made fair deliberately: the same task, the same split, the same
metric, and the same persistence baseline. The results are reported in
Chapter~\ref{ch:testing}; the outcome was that the sequence model is genuinely
competitive at the single horizon it was built for and was nonetheless not
promoted, for reasons of coverage, cost and explainability that
Chapter~\ref{ch:testing} sets out. The module remains in the repository as
evidence, and is excluded from the deployment image so that its dependency never
reaches production.

\section{The Inference Pipeline}
\label{sec:impl-inference}
Inference loads the champion once and then processes each city independently. For
a city it fetches the past three days and the next four days in a single request,
engineers features over the combined series - so that the anchor's rolling
statistics are computed from real history rather than being unavailable - and
splits the result at the current hour. The last row at or before now becomes the
anchor; the first 72 rows after now become the targets.

The model input matrix is then assembled column by column: target-time features
from the future rows, anchor-state features broadcast from the single anchor row,
and a horizon column computed as the whole number of hours between each future row
and now. The matrix is reordered to the bundle's stored column order before
prediction, which is the guard against skew mentioned earlier.
Figure~\ref{fig:code-infer} shows the assembly.

\begin{figure}[htbp]\centering
\begin{lstlisting}[language=Python]
X = pd.DataFrame(index=future.index)
for col in TARGET_TIME_FEATURES:      # forecast weather + calendar at tau
    X[col] = future[col].values
for col in ANCHOR_STATE_FEATURES:     # observed state now, broadcast
    X[f"{col}_anchor"] = anchor[col]
X["horizon"] = ((future["datetime"] - now) / pd.Timedelta("1h")) \
                 .round().astype(int).values
X = X[bundle.feature_columns]         # exact training column order
preds = bundle.estimator.predict(X)
\end{lstlisting}
\codegap
\caption{Assembling 72 model inputs for one city from one anchor row and 72
future rows.}
\label{fig:code-infer}
\end{figure}

Each prediction is then widened by the residual quantiles stored for its horizon,
clipped to the valid index range, and labelled with its health category; the
hourly series is aggregated into three calendar days; the alert module scans the
curve; and the peak hour is explained. Every city is wrapped in an exception
handler, so one city's failure costs one city.

The whole payload is written as a single document and, separately, every forecast
point is appended to the forecast log for later scoring. Both the explanation and
the logging steps are wrapped so that their failure cannot prevent the forecast
from being published, in accordance with principle P5.

\section{Explainability}
\label{sec:impl-shap}
Explanation uses TreeSHAP~\cite{ref:treeshap} against the champion estimator. For
a single input row it computes the exact Shapley values, takes the model's base
value, and reports the five largest contributions by absolute magnitude. Because
the attribution is additive, the base value plus the sum of all contributions is
exactly the prediction, which means the explanation can be phrased arithmetically
rather than as a bare ranking.

Two implementation choices make the output usable by a non-specialist. The
explainer is cached across calls within a process, since constructing it is the
expensive part and inference explains 22 cities in one job. And every feature name
is mapped through a dictionary of human-readable labels, so that
\texttt{aqi\_roll\_mean\_6h\_anchor} is presented as ``6h AQI trend'' and
\texttt{apparent\_temperature} as ``feels-like temp''.
Figure~\ref{fig:code-shap} shows the rendering, and the sentence it produced for
Lahore in the live system is visible in the dashboard screenshot in
Chapter~\ref{ch:ui}: \emph{``Forecast AQI $\approx$ 200 (baseline 132). Main
drivers: current AQI +26.0, time of day +13.4, season $-$6.6, 6h AQI trend +6.4,
temperature +5.0.''}

\begin{figure}[htbp]\centering
\begin{lstlisting}[language=Python]
shap_vals = np.asarray(explainer.shap_values(X))[0]
base = float(np.ravel(explainer.expected_value)[0])
pred = base + float(shap_vals.sum())   # additive: this is exact

order = np.argsort(np.abs(shap_vals))[::-1][:top_k]
contributors = [
    {"label": friendly(bundle.feature_columns[i]),
     "impact": round(float(shap_vals[i]), 1)} for i in order]
parts = [f"{c['label']} {'+' if c['impact'] >= 0 else '-'}"
         f"{abs(c['impact'])}" for c in contributors]
text = (f"Forecast AQI ~ {round(pred)} (baseline {round(base)}). "
        f"Main drivers: " + ", ".join(parts) + ".")
\end{lstlisting}
\codegap
\caption{Rendering Shapley attributions as a plain-language sentence. Additivity
is what allows the numbers to be stated arithmetically.}
\label{fig:code-shap}
\end{figure}

A global view is produced separately by averaging absolute Shapley values over a
sample and plotting the leading features; the resulting ranking is reported in
Chapter~\ref{ch:testing}.

\section{The What-If Simulator}
\label{sec:impl-whatif}
The simulator exposes five drivers: current particulate concentration, current
index, wind speed, humidity and temperature. It first builds the model input for
the selected city at a lead time of 24 hours from live data, and returns the
current real value of each driver, so that the interface's sliders begin at
reality rather than at an arbitrary midpoint. The user's overrides are then mapped
onto the underlying feature columns and the model is called twice, once
unmodified and once overridden, and both values are returned with their health
categories and the difference between them.

The mapping deserves comment because it is where the simulator earns its
plausibility. Overriding the current index sets not only the anchor index but also
the 6-hour and 24-hour anchor means, since a stated current level that contradicts
its own recent history would place the model far outside its training
distribution. Overriding the wind speed rescales both wind vector components so
that their resultant matches the requested speed while preserving direction, and
adjusts the gust feature proportionally. Overriding the temperature adjusts the
apparent temperature with it. Each of these keeps the perturbed row internally
consistent, which is what makes the response interpretable.

The simulator is labelled in the interface as a model simulation rather than
causal evidence, in accordance with NFR-6. What it shows is how \emph{this model}
responds to a change in its inputs, which is a genuine and useful statement about
the model, and not a claim about what would physically happen.

\section{Monitoring}
\label{sec:impl-monitoring}
Monitoring implements the design of Chapter~\ref{ch:design} in two halves, shown
in Figure~\ref{fig:monitoring}.

\fig{monitoring}{1.0}{The two monitoring signals and why they are read
together.}{fig:monitoring}

The drift half divides the feature history at fourteen days before the latest
observation, treats everything older as the reference and everything newer as the
current window, and computes the Population Stability Index per monitored feature.
The implementation, in Figure~\ref{fig:code-psi}, takes the bin edges from
deciles of the reference distribution, deduplicates them so that a degenerate
distribution cannot produce empty bins, replaces the outer edges with infinities
so that new extreme values fall inside the outermost bins rather than being
dropped, and adds a small constant to every proportion so that an empty bin
cannot produce an infinite logarithm. It returns zero rather than a spurious
value when either window is too small to support the estimate.

\begin{figure}[htbp]\centering
\begin{lstlisting}[language=Python]
def population_stability_index(ref, cur, bins: int = 10) -> float:
    """PSI: <0.1 stable, 0.1-0.25 moderate, >0.25 significant."""
    ref, cur = ref.dropna(), cur.dropna()
    if len(ref) < 50 or len(cur) < 20:
        return 0.0
    edges = np.unique(np.quantile(ref, np.linspace(0, 1, bins + 1)))
    if len(edges) < 3:
        return 0.0
    edges[0], edges[-1] = -np.inf, np.inf   # absorb unseen extremes
    ref_pct = np.histogram(ref, edges)[0] / len(ref) + 1e-6
    cur_pct = np.histogram(cur, edges)[0] / len(cur) + 1e-6
    return float(np.sum((cur_pct - ref_pct)
                        * np.log(cur_pct / ref_pct)))
\end{lstlisting}
\codegap
\caption{The Population Stability Index, with the guards that keep it defined on
real data.}
\label{fig:code-psi}
\end{figure}

The error half appends every issued forecast to a columnar log, deduplicated
on the triple of issue time, city and target hour, and later joins that log to the
observed index on city and target hour. Only rows for which reality has arrived
can be scored, so the report distinguishes between having no forecasts logged,
having forecasts logged but no matured outcomes, and having a real score -
rather than reporting a misleading zero in any of those cases.

Both halves are written to a committed document each night, which is what allows
the deployed backend to serve the monitoring tab without any live feature-store
connection.

\section{Alerting}
The alerting module scans a city's hourly forecast and produces a graded alert.
Two thresholds are used, aligned with the EPA category boundaries: 150, the start
of the \emph{unhealthy} band, raises a warning, and 200, the start of \emph{very
unhealthy}, raises a danger. The alert records the peak value, the hour of the
peak, the category and its standard advice, and - the field that makes it
actionable - the first hour at which the curve crosses the warning threshold, so
a user is told when the episode begins rather than only how bad it will get.

\section{The API}
The backend exposes the endpoints listed in Table~\ref{tab:api}. Its structure
follows the design rule that document reads and model invocations are separate.
Document endpoints read a committed file and are cheap; the model is loaded lazily
and only by the endpoints that genuinely need it, so a deployment that never
receives a what-if request never loads the estimator at all.

The monitoring endpoint illustrates the layered error handling of principle P5 in
its final form. It first serves the committed snapshot; failing that it attempts
to compute the report live; failing that it returns a well-formed placeholder
explaining that monitoring runs in the training pipeline and a snapshot will
appear after the next run. The tab therefore reports its own state honestly and
never presents an error to a member of the public.

Cross-origin requests are permitted because the frontend is served from a
different origin and every endpoint is a read of public data. The advisor endpoint
checks for its credential before attempting anything and returns a specific,
actionable message when it is absent, rather than failing obscurely inside the
client library.

\section{The Advisor and the Tool Server}
The advisor answers natural-language questions by tool use rather than by recall.
Four functions are declared to the model with JSON schemas: list the supported
cities, get a city's current index and three-day forecast with its alert, get a
city's long-run historical statistics, and explain why a city's forecast is high
or low. The model is invited to call them; the backend executes the call against
the real forecast document and returns the result; the loop repeats up to a
bounded number of rounds until the model produces a final answer. The system
prompt instructs the model to ground every number in a tool call and to answer
with practical health guidance - who is at risk, whether to limit outdoor
exertion, whether a mask is warranted, and when the cleanest part of the day is.

The same four functions are published over the Model Context
Protocol~\cite{ref:mcp} by a small server module, so that an external protocol
client can query Pakistani air quality through exactly the same code path that
serves the dashboard's advisor. Sharing the implementation rather than duplicating
it is what guarantees the two cannot disagree.

\section{Continuous Integration and Scheduling}
\label{sec:impl-cicd}
Three workflows constitute the compute plane.

The \textbf{feature workflow} runs hourly at five minutes past, checks out the
repository, installs the package, and runs the feature pipeline with the
feature-store credentials injected from the secret store.

The \textbf{training workflow} runs daily at 02:30~UTC and executes four steps in
order: train, infer, evaluate and monitor, then commit the refreshed artefacts
back to the repository. The evaluation and monitoring steps are permitted to fail
without failing the workflow, since neither is required for a forecast to be
served. The commit message carries a skip marker so that committing artefacts
cannot trigger another run.

The \textbf{backfill workflow} is manual, takes the history start date as an
input, and is given a much longer time allowance.

Two mechanisms keep these safe. The feature and backfill workflows share a
concurrency group, so the scheduled hourly job and a long manual backfill can
never write to the feature store simultaneously; and every workflow has an
explicit timeout, so a stalled job cannot consume the free allowance indefinitely.

The commit step is what closes the loop between the compute plane and the serving
layer: the refreshed forecast, evaluation, monitoring and leaderboard documents
become part of the repository, and are therefore present in the next deployment
image and readable by the static frontend.

\section{Deployment}
\label{sec:impl-deploy}
The deployment topology was shown in Figure~\ref{fig:deployment}.

The backend is a container built from a slim Python base. Three details of the
image were arrived at empirically. The OpenMP runtime is installed explicitly,
because the gradient-boosting library fails at import time without it on a slim
image. The dependency set is a reduced one that deliberately omits the
feature-store client and the deep-learning framework, because the deployed API
serves committed documents and a bundled model and needs neither; this reduces the
image substantially. And the package is installed in editable mode, so that the
project-root path derived in the configuration module resolves to the application
directory where the committed documents and the model bundle actually live.

The frontend is a static build served from an edge host. The backend's address is
supplied by baking it into the build command rather than by setting a build-time
environment variable, because the hosting platform rejects variables carrying the
build tool's public prefix as potentially sensitive. Since the address is a public
URL, baking it in is appropriate rather than a workaround. A rewrite rule directs
every path to the application, so that the single-page router works on a direct
navigation.

\section{Challenges Encountered and Their Resolution}
\label{sec:impl-challenges}
Table~\ref{tab:challenges} records the substantive failures encountered during
development, their root causes and their resolutions. Several of them shaped the
design rather than merely being fixed, and they are the most transferable findings
in this thesis, because they are what the free-tier, serverless constraint
actually costs.

\begin{xltabular}{\textwidth}{P{3.4cm} P{4.0cm} Y}
\caption{Challenges encountered and their resolution.}\label{tab:challenges}\\
\toprule
\bfseries Challenge & \bfseries Root cause & \bfseries Resolution \\
\midrule
\endfirsthead
\toprule
\bfseries Challenge & \bfseries Root cause & \bfseries Resolution \\
\midrule
\endhead
\bottomrule
\endlastfoot
Feature group could not be created & The platform's newer default table format
requires a client library the base installation does not ship. & Created the
group with the server-side materialised time-travel format, which preserves
primary-key upserts. \\
Materialisation job failed with an out-of-memory error & The statistics profiler
ran over all 74 columns on a free-tier executor. & Disabled statistics on the
feature group and created a fresh version; statistics are used only by the
platform's interface and by no pipeline. \\
Backfill froze at the tenth of 22 cities & A blocking insert waited indefinitely
on a materialisation job stalled at initialisation, until the scheduler's time cap
killed the run. & Made inserts non-blocking and wrote city by city; safe because
writes are idempotent upserts and the resume is coverage-based. \\
Backfill skipped cities that were not in fact loaded & The resume checked whether
a city was \emph{present}, but the hourly pipeline had already seeded every city
with recent rows. & Changed the resume to check whether a city's earliest stored
observation reaches back to the requested start date. \\
Training crashed in the chronological split & The feature store returns the event
time column as strings, so a quantile over it was undefined. & Coerced to a
datetime type on read and again inside the split, so the split is correct
regardless of how the store returned the column. \\
The champion and the leaderboard reset on every run & Both were being written to
the runner's local directory, which is destroyed when the job ends. & Pushed the
champion to the durable model registry and committed the leaderboard to the
repository. \\
The feature-store client could not be installed for local development & A
transitive cryptographic dependency fails to build on the development platform. &
Introduced the storage facade with an interchangeable local columnar backend;
all feature-store work runs on Linux automation. \\
The container build failed to start & The platform's default builder ignored the
declared build configuration and no start command was inferred. & Added a
Dockerfile, which the platform detects in preference and which gives explicit
control of the base image, the system packages and the entry point. \\
The deployed backend reported that no forecast was available & The deployment
tool honours the repository's ignore file, so the committed model bundle and data
documents were excluded from the uploaded context. & Added explicit negations for
those paths to the ignore file so that the required artefacts are shipped. \\
The frontend could not be given the backend address & The hosting platform rejects
build-time variables carrying the build tool's public prefix. & Baked the address
into the build command in the platform configuration; the address is public, so
this is appropriate rather than a compromise. \\
The gradient-boosting library failed at import in the container & The slim base
image lacks the OpenMP runtime. & Installed it explicitly in the image. \\
The application resolved its data paths to the wrong directory in the container &
An ordinary install placed the package in the interpreter's library directory, so
the derived project root pointed away from the committed artefacts. & Installed
the package in editable mode inside the image. \\
\end{xltabular}

\section{Summary of Implementation}
This chapter described the construction of the system: the technology stack and
its justification; the single configuration module; the retrying ingestion client;
the EPA index implementation and its deliberate restriction to particulate
species; the feature builder and the two mechanisms - shifted windows and
per-city grouping - that prevent leakage; the multi-horizon supervised builder;
the hourly, backfill, training and inference pipelines; the evaluation suite; the
champion-challenger gate and the durable registry, including the refused tie that
demonstrates the gate working; the model layer and the sequence-model study;
explanation, simulation, monitoring and alerting; the API and its layered
fallbacks; the advisor and the tool server; the three scheduled workflows; the
container and static deployments; and the twelve substantive failures encountered
along the way together with what each of them required. The next chapter presents
the user interface through which all of this reaches a member of the public.
